% !TEX root = ../thesis.tex

\chaptermark{Predhovor}
\phantomsection{}
\addcontentsline{toc}{chapter}{Predhovor}

\chapter*{Predhovor}\label{ch:preface}

To, ako som sa dostal k~počítačom, si pamätám veľmi dobre dodnes. Bolo to ešte na základnej škole, keď som ešte ako \uv{malý jogurt} v~treťom ročníku základnej školy (takže som mal asi 8 alebo 9 rokov) navštívil počas školy v~prírode \emph{Dom pionierov} v~Trstenej. Posadili nás pred 8-bitové mikropočítače \emph{PMD 85-2A}\footnote{\url{https://sk.wikipedia.org/wiki/PMD_85}}, ktoré boli pripojené k~televízorom \emph{Merkúr} a na nich svietil mikrosvet robota \emph{Karla}, ktorého sme v~to poobedie ovládali. Už vtedy som vedel, že keď budem veľký, tak už nechcem byť stolárom, ale programátorom.

Odvtedy už prešlo veľa rokov. Medzičasom som robota \emph{Karla} postavil ako papierový model, čítal som o~ňom v~časopise \emph{Zenit pionierov} (neskôr \emph{Elektrón-Zenit}), ale už som v~ňom neprogramoval. Programovať som sa naučil v~jazyku \emph{BASIC} na domácom mikropočítači \emph{Didaktik M} (domáci klon populárneho mikropočítača \emph{ZX Spectrum}), a keďže som rástol v~90-tych rokoch, postupne som z~8-bitov prešiel na 32-bitovú architektúru zastúpenú osobnými počítačmi označovanými ako \emph{PC}.

Ku robotovi \emph{Karlovi} som sa opäť dostal až o~mnoho rokov neskôr, keď sme na \emph{Katedre počítačov a informatiky} \emph{Fakulty elektrotechniky a informatiky} \emph{Technickej univerzity v~Košiciach} rozmýšľali, ako inovovať štúdium úvodných kurzov programovania v~programovacom jazyku \emph{C} s~\emph{Pecinovského} \uv{umožnite študentom programovať čo najskôr} \cite{pecinovsky}. Pri hľadaní inšpirácií, ako vyzerajú úvodné kurzy programovania na iných, najmä svetových univerzitách, som sa dostal na \emph{Stanfordskú univezitu} do kurzu \emph{Programming Methodology (CS106A)}\footnote{\url{https://see.stanford.edu/Course/CS106A}}. A~práve v~tomto kurze a práve na tejto univezite síce začínajú s~programovaním v~jazyku \emph{Java}, ale práve pomocou sveta robota \emph{Karla}.

Použiť robota \emph{Karla} v~našich úvodných kurzoch programovania v~jazyku \emph{C} pre našich študentov v~prvom ročníku mi prišiel ako veľmi dobrý nápad. V~jeho použití som videl obrovskú výhodu v~tom, že dokážeme so študentmi vytvárať programy okamžite od prvého cvičenia. Práve to nie je až tak jednoduché v~prípade jazyka \emph{C} a to nie len kvôli operátorm referencie a dereferencie.

Problém však nastal pri hľadaní vhodnej knižnice, ktorá by robota \emph{Karla} v~jazyku \emph{C} poskytovala. Aj napriek tomu, že existujú rozličné hlavne grafické implementácie s~vlastným pseudojazykom, môj zámer je iný. Chcem knižnicu, pomocou ktorej by bolo možné \emph{Karla} ovládať pomocou samostatných funkcií, pričom vizualizácia by bola v~textovom (konzolovom) rozhraní. A~dosiahnutiu tohto cieľa je venovaná táto práca.


% Úvod práce stručne opisuje stanovený problém, kontext problému a motiváciu pre riešenie problému. Z~úvodu by malo byť jasné, že stanovený problém doposiaľ nie je vyriešený a má zmysel ho riešiť.
% V~úvode neuvádzajte štruktúru práce, t.j. o~čom je ktorá kapitola. Rozsah úvodu je minimálne 2 celé strany (vrátane formulácie úlohy).

